\subsection{交的性质}

引理 7. 设 E 是右 A-模，F 是左 A-模， \(\mathbf{F}'\) 、 \(\mathbf{F}''\) 是 F 的两个子模，使得 \(\mathbf{F} = \mathbf{F}' + \mathbf{F}''\) 。则 \(\mathbf{E} \otimes \mathbf{F}'\) 和 \(\mathbf{E} \otimes \mathbf{F}''\) 在 \(\mathbf{E} \otimes \mathbf{F}\) 中的典范像的交等于 \(\mathbf{E} \otimes (\mathbf{F}' \cap \mathbf{F}'')\) 的典范像。

考虑图

\[
\begin{array}{c} 0 \longrightarrow \mathrm{F} ^ {\prime} \cap \mathrm{F} ^ {\prime \prime} \longrightarrow \mathrm{F} ^ {\prime} \longrightarrow \mathrm{F} ^ {\prime} / (\mathrm{F} ^ {\prime} \cap \mathrm{F} ^ {\prime \prime}) \longrightarrow 0 \\ \Biggl \downarrow \quad \Biggl \downarrow \quad \text {   j   } \\ 0 \longrightarrow \mathrm{F} ^ {\prime \prime} \longrightarrow \mathrm{F} ^ {\prime} + \mathrm{F} ^ {\prime \prime} \longrightarrow (\mathrm{F} ^ {\prime} + \mathrm{F} ^ {\prime \prime}) / \mathrm{F} ^ {\prime \prime} \longrightarrow 0 \end{array}
\]

其中未标明的箭头是典范单射和满射，j 是《代数学》第一章 §6 第 13 目定理 6 中定义的典范同构。这个图是交换的且它的行是正合的。我们（由于 \(F = F' + F''\) ）得到一个交换图

\[
\begin{array}{c} \mathbf {E} \otimes (\mathbf {F} ^ {\prime} \cap \mathbf {F} ^ {\prime \prime}) \longrightarrow \mathbf {E} \otimes \mathbf {F} ^ {\prime} \longrightarrow \mathbf {E} \otimes (\mathbf {F} ^ {\prime} / (\mathbf {F} ^ {\prime} \cap \mathbf {F} ^ {\prime \prime})) \\ \Biggl \downarrow \qquad \qquad \qquad \qquad \qquad \qquad \qquad \qquad \qquad \qquad \qquad \qquad \qquad 1 _ {\mathbf {E}} \otimes \mathbf {j} \\ \mathbf {E} \otimes \mathbf {F} ^ {\prime \prime} \longrightarrow \mathbf {E} \otimes \mathbf {F} \longrightarrow \mathbf {E} \otimes (\mathbf {F} / \mathbf {F} ^ {\prime \prime}) \end{array}
\]

这个图的行是正合的（第 1 目，引理 1），且 \(1_{E} \otimes j\) 是同构。于是我们的断言是 §1 第 4 目命题 1 (i) 的特殊情形（参看习题 5）。

命题 6. 设 E 是右 A-模，F 是左 A-模，且 E 对 F 是平坦的。对 F 的每个子模 \(\mathbf{F}'\) ，我们用 \(\phi(\mathbf{F}')\) 表示 \(\mathbf{E} \otimes \mathbf{F}'\) 在 \(\mathbf{E} \otimes \mathbf{F}'\) 到 \(\mathbf{E} \otimes \mathbf{F}\) 的典范映射（由第 2 目的定义 1，它是单射的）之下的像。则若 \(\mathbf{F}'\) 、 \(\mathbf{F}''\) 是 F 的两个子模，我们有

\[
\phi (\mathbf {F} ^ {\prime} \cap \mathbf {F} ^ {\prime \prime}) = \phi (\mathbf {F} ^ {\prime}) \cap \phi (\mathbf {F} ^ {\prime \prime}).
\]

由于 E 对 F 平坦， \(\phi(F' + F'')\) 等同于 \(E \otimes (F' + F'')\) ，而子模 \(\phi(F')\) 、 \(\phi(F'')\) 和 \(\phi(F' \cap F'')\) 分别等同于 \(E \otimes F'\) 、 \(E \otimes F''\) 和 \(E \otimes (F' \cap F'')\) 在 \(E \otimes (F' + F'')\) 中的典范像。于是命题 6 由引理 7 得出。

注记 1. 在命题 6 的假设下，通常把 \(\mathbf{E} \otimes \mathbf{F}'\) 等同于 \(\phi(\mathbf{F}')\) （对 \(\mathbf{F}'\) 是 \(\mathbf{F}\) 的每个子模），由此得到公式

\[
\mathrm{E} \otimes_ {\mathrm{A}} (\mathrm{F} ^ {\prime} \cap \mathrm{F} ^ {\prime \prime}) = (\mathrm{E} \otimes_ {\mathrm{A}} \mathrm{F} ^ {\prime}) \cap (\mathrm{E} \otimes_ {\mathrm{A}} \mathrm{F} ^ {\prime \prime}).
\]

命题 7. 设 E 是右 A-模，E‘ 是 E 的子模，F 是左 A-模，F’ 是 F 的子模。设 E/E‘ 或 F/F’ 是平坦的。则 \(E' \otimes F'\) 在 \(E \otimes F\) 中的典范像等于 \(E' \otimes F\) 和 \(E \otimes F'\) 在 \(E \otimes F\) 中的典范像的交。

例如设 \(E/E'\) 是平坦的，考虑图

\[
\begin{array}{c} \mathrm{E} ^ {\prime} \otimes \mathrm{F} ^ {\prime} \longrightarrow \mathrm{E} \otimes \mathrm{F} ^ {\prime} \longrightarrow (\mathrm{E} / \mathrm{E} ^ {\prime}) \otimes \mathrm{F} ^ {\prime} \\ \Biggl \downarrow \qquad \qquad \qquad \qquad \qquad \qquad \qquad \qquad \qquad \qquad \qquad \qquad \qquad u \Biggl \downarrow \\ \mathrm{E} ^ {\prime} \otimes \mathrm{F} \longrightarrow \mathrm{E} \otimes \mathrm{F} \longrightarrow (\mathrm{E} / \mathrm{E} ^ {\prime}) \otimes \mathrm{F} \end{array}
\]

其中箭头是典范同态。这个图是交换的且它的行是正合的（第 1 目，引理 1）。由于 E/E' 是平坦的，u 是单射的。于是我们的断言是 §1 第 4 目命题 1 (i) 的特殊情形。

推论. 设 E 是右 A-模，E' 是 E 的子模。

\begin{enumerate}
\def\labelenumi{(\roman{enumi})}
\tightlist
\item
  设 \(\mathbf{E} / \mathbf{E}'\) 是平坦的。则对 \(\mathfrak{a}\) 取 \(\mathbf{A}\) 的每个左理想，
\end{enumerate}

\[
\mathbf {E} ^ {\prime} \mathfrak {a} = \mathbf {E} ^ {\prime} \cap \mathbf {E} \mathfrak {a}.\tag{5}
\]

\begin{enumerate}
\def\labelenumi{(\roman{enumi})}
\setcounter{enumi}{1}
\item
  反之，设 E 是平坦的，且对 A 的每个有限生成左理想 a，关系式 (5) 成立。则 E/E' 是平坦的。
\item
  只需把命题应用于 \(\mathbf{F} = \mathbf{A}_s\) ， \(\mathbf{F}' = a\) 的情形。
\item
  为证明 \(E/E'\) 是平坦的，应用第 5 目命题 4 的判别法 (c)；这时需要证明序列
\end{enumerate}

\[
0 \rightarrow \mathrm{E} ^ {\prime} / \mathrm{E} ^ {\prime} a \rightarrow \mathrm{E} / \mathrm{E} a \rightarrow \mathrm{E} / (\mathrm{E} ^ {\prime} + \mathrm{E} a) \rightarrow 0
\]

对 A 的每个有限生成左理想 a 在 E'/E'a 处正合。而这正是关系式 (5) 所表达的。

注记 2. 如果只假设 \(\mathbf{E} / \mathbf{E}'\) 对 F 平坦或 \(\mathbf{F} / \mathbf{F}'\) 对 E 平坦，命题 7 的结论仍然成立。

